\chapter{Introduction}
\pagenumbering{arabic}
\vspace{-25pt}

For the remainder of this chapter: let $U \subset \bfR^n$ be an open set. An element $x \in \bfR^n$ will always be of the form $(x_1,\ldots,x_{n})$. The function $u:U \rightarrow \bfR$ will always denote the unknown function being solved for in a differential equation. We denote $u := u(x) = u(x_1,\ldots,x_{n})$.


\section{Basic Definitions}

\begin{definition}
	The \textui{unit vector in $\bfR^n$} is denoted by $e^i$, where it is $1$ is the $i^\text{th}$ slot and $0$ otherwise.
\end{definition}

\begin{definition}
	The \textui{partial derivative} of $u$ with respect to $x_i$ is:
	\[
		u_{x_i} := \frac{\partial u}{\partial x_i} := \lim_{h \rightarrow 0} \frac{u(x+he^i) - u(x)}{h}.
	\] 
\end{definition}

\begin{definition}
	The \textui{gradient} of $u$ is the vector:
	\[
		Du := \nabla u := (u_{x_1},\ldots,u_{x_n}).
	\] 
\end{definition}

\begin{definition}
	The \textui{Hessian matrix}  of $u$ is the 2 dimensional array:
	\[
		D^{2}u = \bmat u_{x_1 x_1} & ... & u_{x_1 x_{n}} \\ \vdots & \ddots & \vdots \\ u_{x_{n} x_1} & ... & u_{x_{n} x_{n}} \emat
	\] 
\end{definition}

\begin{definition}
	The \textui{Laplacian} of $u$ is:
	\[
		\Delta u := \nabla^2 u := \Tr(D^2 u) = \sum_{i=1}^{n} u_{x_i x_i}.
	\] 
\end{definition}

\begin{definition}
	\phantom{a}
	\begin{enumerate}[label = (\alph*),itemsep=1pt,topsep=3pt]
		\item A vector of the form $\alpha := (\alpha_1,\ldots,\alpha_{n})$, where $\alpha_i \in \bfZ_{\geq 0}$, is called a \textui{multiindex of order $|\alpha|$}, where:
			\[
				|\alpha| = \alpha_1 + \ldots +\alpha_n.
			\]

		\item Give a multiindex $\alpha$, we define:
			\[
				D^\alpha u(x) := \frac{\partial^{|\alpha|}u(x)}{\partial^{\alpha_1}x_1\ldots\partial^{\alpha_n}x_n}.
			\] 
	\end{enumerate}
\end{definition}

\begin{definition}
	A \textui{partial differential equation} (PDE for short) is an equation involving an unknown function $u$ along with one or more of its partial derivatives. A partial differential equation is  \textui{linear} if we can write it as:
	\[
		\sum_{|a| \leq K}^{} a_{\alpha}(x)D^{\alpha}u(x) = f(x).
	\]
	If $f(x)=0$, then this PDE is called \textui{homogenous}. 
\end{definition}

\section{Multivariable Calculus}

\begin{theorem}[Fubini]
	If $u \in L^1(U)$, then:
	\[
		\int_{U}^{} u dx = \underbrace{\int \ldots \int}_{n\text{-times}}u(x_1,\ldots,x_{n})dx_1\ldots dx_n.
	\] 
\end{theorem}

\begin{definition}
	The \textui{boundary} of $U$ is denoted $\partial U$. 
\end{definition}

\begin{definition}
	We say the boundary $\partial U$ is \textui{$C^k$} if in a neighborhood of any point on the boundary, the boundary is the graph of a $C^k$ function (possibly after the coordinate system has been rotated).
\end{definition}

\begin{definition}
	If $\partial U$ is $C^1$, then along $\partial U$ is defined the \textui{outward pointing unit normal vector field}, which we denote by:
	\[
		\vec{\nu} = (\nu^1,\ldots,\nu^n).
	\] 
\end{definition}

\begin{definition}
	Let $u \in C^1(\overline{U})$. The \textui{outward normal derivative} of $u$ is:
	\[
		\frac{\partial u}{\partial \nu} := Du \cdot \vec{\nu}.
	\] 
\end{definition}

\begin{definition}
	Let $\vec{u} = (u^1,\ldots,u^n)$ be a vector field. The \textui{flux of $u$ across $\partial U$} is:	
	\[
		\int_{\partial U}\vec{u}\cdot \vec{\nu}dS. 
	\] 
\end{definition}

\begin{theorem}[Gauss-Green]
	\phantom{a}
	\begin{enumerate}[label = (\alph*),itemsep=1pt,topsep=3pt]
		\item Suppose $u \in C^1(\overline{U})$. For any $i=1,\ldots,n$, we have:
		\[
			\int_{U}^{} u_{x_i}dx = \int_{\partial U}^{} u \nu^i dS
		\]
		\item (Divergence Theorem) Let $\vec{v} = (v^1,\ldots,v^n)$, where $v^i \in C^1(\overline{U})$ for each $i=1,\ldots,n$. Then:
			\[
				\int_{U}^{} \Div \vec{v}dx = \int_{\partial U}^{} \vec{v}\cdot \vec{\nu}dS.
			\] 	
	\end{enumerate}
\end{theorem}

\begin{theorem}[Integration by Parts]
	Let $u,v \in C^1(\overline{U})$ and let $\partial U$ be $C^1$. Then:
	\[
		\int_{U}^{} u_{x_i}v dx = - \int_{U}^{} u v_{x_i} dx + \int_{\partial U}^{} uv \nu^i dS.
	\] 
\end{theorem}

\begin{theorem}[Green's Formulas]
	Let $u,v \in C^1(\overline{U})$ and let $\partial U$ be $C^1$. Then:
	\begin{enumerate}[label = (\alph*),itemsep=1pt,topsep=3pt]
        \item $\ds \int_U \Delta u dx = \int_{\partial U}\frac{\partial u}{\partial \nu}dS$,
        \item $\ds \int_U Dv \cdot Du dx = -\int_U u\Delta v dx + \int_{\partial U}\frac{\partial u}{\partial \nu} u dS$, 
        \item $\ds \int_U u \Delta v - v \Delta u dx = \int_{\partial U} u\frac{\partial v}{\partial \nu} - v \frac{\partial u}{\partial \nu}dS$.
    \end{enumerate}
\end{theorem}
	
\section{Polar Coordinated in $\bfR^n$ }

\begin{remark}
	In $\bfR^n$, for any $x \neq 0$, we can write:
	\[
	x = |x| \cdot \frac{x}{|x|} := r \omega,
	\] 
	where $r\in (0,\infty)$ and $\omega \in S^{n-1} = \partial B(0,1)$.
\end{remark}

\begin{theorem}
	There exists a measure $d \omega$ on $S^{n-1}$ such that:
	\[
		\int_{\bfR^n}^{}u dx = \int_{0}^{\infty} \left( \int_{S^{n-1}}^{} u(r \omega) r^{n-1}d \omega \right) dr.
	\] 
	In other words, $dx = r^{n-1}d\omega dr$.
\end{theorem}

\begin{corollary}
	For any $x_0 \in \bfR^n$ and $a > 0$, we have:
	\[
		\int_{B(x_0,a)}^{} u dx = \int_{B(0,a)}^{} u(x_0 + y)dy = \int_{0}^{a} \left( \int_{S^{n-1}}^{} u(x_0+r \omega)r^{n-1} d\omega \right) dr.
	\] 
\end{corollary}

\begin{corollary}
	For any  $x_0 \in \bfR^n$ and $r > 0$, we have:
	\[
	\frac{d}{dr}\int_{B(x_0,r)}^{} u dx = \int_{\partial B(x_0,r)}^{} udS.
	\] 
\end{corollary}

\section{Smoothing}

\begin{definition}
	\phantom{a}
	\begin{enumerate}[label = (\alph*),itemsep=1pt,topsep=3pt]
		\item Let $U,V \subset \bfR^n$. We say $U$ is \textui{compactly contained} in $V$, and write $U \subset\subset V$, if $V \subset \overline{V} \subset U$ and $\overline{V}$ is compact. 
		\item If $\restr{f}{V} \in L^p(V)$ for all $V \subset\subset U$, then we say $f$ is \textui{locally integrable on $U$}. The set of all such functions is denoted $L_{\text{loc}}^p(U)$.
	\end{enumerate}
\end{definition}

\begin{definition}
	For $\epsilon > 0$, we define:
	\[
		U_\epsilon := \{x \in U \mid \dist(x, \partial U) > \epsilon\}.
	\] 
\end{definition}

\begin{definition}
	\phantom{a}
	\begin{enumerate}[label = (\alph*),itemsep=1pt,topsep=3pt]
		\item Let $x \in \bfR^n$. The \textui{standard mollifier} is the function:
		\[
			\eta(x) := 
				\begin{cases}
			 		\frac{-C}{1-|x|^2}, & |x| < 1 \\
					0, & |x| \geq 1
				\end{cases},
		\] 
		where $C$ is chosen so that $\int_{\bfR^n}^{} \eta(x)dx = 1$.

		\item For all $\epsilon > 0$, define:
			\[
				\eta_\epsilon(x) := \frac{1}{\epsilon^n}\eta\left( \frac{x}{\epsilon}\right).
			\] 
	\end{enumerate}
\end{definition}

\begin{proposition}
	We have $\eta \in C^\infty_c(\bfR^n)$. Moreover, for all $\epsilon > 0$, we have:
	\begin{enumerate}[label = (\alph*),itemsep=1pt,topsep=3pt]
		\item $\eta_\epsilon \in C^\infty_c(\bfR^n)$;
		\item $\supp \eta_\epsilon = B(0,\epsilon)$;
		\item $\int_{\bfR^n}^{} \eta_\epsilon(x)dx = 1$.
	\end{enumerate}
\end{proposition}

\begin{definition}
	\phantom{a}
	\begin{enumerate}[label = (\alph*),itemsep=1pt,topsep=3pt]
		\item For $\epsilon > 0$, we define:
			\[
				U_\epsilon := \{x \in U \mid \dist(x, \partial U) > \epsilon\}.
			\] 
		\item If $f \in \Loc^p(U)$, we define it \textui{mollifier} as:
			\begin{align*}
				f^{\epsilon}(x)
				&:=   \restr{(\eta_\epsilon \ast f)}{U_\epsilon}\\
				&= \int_{U}^{} \eta_\epsilon(x-y)f(y)dy \\
				&= \int_{B(0,\epsilon}^{} \eta_\epsilon(y)f(x-y)dy. \\
			\end{align*}
		
	\end{enumerate}
\end{definition}




